\section{Walk expansion without scalar trace masses}\label{s:walks}

We work on a finite subset of the plane with the laws of
Definition~\ref{d:moment}. We first prove the finite expansion and
uniform response estimates, and then give the local and open-trail
contracts used to bound it.

\begin{definition}\label{d:domain}
Fix finite \(G\), \(S=S(G)\), and original unit pins \(\tau\).
A domain assigns statuses \(\mathrm U,\mathrm D,\mathrm X\);
\(\mathrm X\) sites are removed. Each nonremoved site has a kept
subset of its incident bonds inside \(G\). A bond is present if both
ends are nonremoved and both keep it. A \(\mathrm D\) site has at
most two present bonds. Protected sites in \(S\) have status
\(\mathrm U\) and keep every incident bond inside \(G\).
Use \(\lambda_{\mathrm U},\lambda_{\mathrm D}\) at the corresponding
sites and define
\[
 Z_D^\tau[f]=\sum_n f(n)
       \prod_{\text{present }\{x,y\}}(1+n_x\cdot n_y)
       \prod_{\substack{(x,z)\in\mathcal L(G)\\x\in S}}
                   (1+\tau_{xz}\cdot n_x)
       \prod_{x:\,\mathrm{st}(x)\ne\mathrm X}\lambda_{\mathrm{st}(x)}(n_x).
\]
The sum is over the product of the supports \(D\) at \(\mathrm U\)
sites and \(D\cup\{0\}\) at \(\mathrm D\) sites.
Define the free sum, normalized expectations and \(\Psi_D\) as before.
The class \(\mathcal D_{\mathrm{cons}}(G)\) consists of domains in
which every \(\mathrm U\) site keeps all its bonds, every \(\mathrm D\)
site keeps exactly two, and every bond between nonremoved sites is
kept by both ends or by neither.
An auxiliary network has the same data but permits protected
\(\mathrm U\) sites to have arbitrary kept subsets. It retains every
original leg factor. Neither domains nor auxiliary networks carry
a scalar trace mass. Only the networks occurring below, with at most
four factors at each site, are used.
\end{definition}

All normalizers in this definition are positive by
Lemma~\ref{l:finitefacts}. Original legs occur only at protected
\(\mathrm U\) sites. In particular every departure site outside
\(S\) is unpinned.
The domains \(G\setminus Y\), \(Y\subset G\setminus S\), in
Definition~\ref{d:pinmodel} are members of
\(\mathcal D_{\mathrm{cons}}(G)\), with statuses
\(\mathrm U,\mathrm X\) only.
An auxiliary network allows a stopped arrival to delete a protected
site's arrival bond while retaining all of that site's original legs.

\begin{definition}\label{d:trace}
A trail uses distinct present unoriented bonds. Its prefixes are
\(w_j=(v_0,\ldots,v_j)\), and \(\mathcal T_D(u,T)\) denotes trails
starting at \(u\) and stopping at their first hit of \(T\).
No trail departs from \(S\).
At departure from \(y=v_{j-1}\), drop \(\{y,v_j\}\) at both ends.
If \(y\) is \(\mathrm U\), change it to \(\mathrm D\), keeping its
other kept bonds. If \(y\) is \(\mathrm D\), remove it.
Call the resulting network \(D[w_j]\), with \(D[w_0]=D\).
Admissible initial starts are unpinned sites outside \(S\) of status
\(\mathrm U\) with at most three present bonds, or status \(\mathrm D\)
with at most one. The same restrictions then hold at every later
departure. Set
\[
 a_D(w)=3^{-N_{\mathrm U}(w)}5^{-N_{\mathrm D}(w)},\qquad
 R_D^\bullet(w)=Z^\bullet_{D[w]}/Z^\bullet_D,
\]
where the counts record the status at departure, not the original status.
Original legs are always unchanged, including at a terminal hit of \(S\).
\end{definition}

We record why this tracing rule is defined on every trail under
consideration. At length zero the stated bounds hold by the start
hypothesis. After an arrival, a \(\mathrm U\) site has lost the
arrival bond from at most four present bonds, so it has at most three;
a \(\mathrm D\) site has lost it from at most two, so it has at
most one. At a \(\mathrm U\) departure, deleting the out-bond
leaves at most two present bonds at the new \(\mathrm D\) site.
At a \(\mathrm D\) departure there is at most the out-bond:
deleting it and removing the site removes no other present bond.
In particular no unused present bond is destroyed by removal.
Induction on the prefix length now shows that every unused bond
of the trail is still present when it is needed, and that every
departure satisfies the asserted bounds.
Sites in \(S\) are never departed from, so none is removed.
They retain all original leg factors even when the final arrival
drops a bond at such a site.

There are at most two departures from any site during this trace:
after its first \(\mathrm U\) departure it is \(\mathrm D\), and
after a \(\mathrm D\) departure it is removed. Also
\[
 |w|=N_{\mathrm U}(w)+N_{\mathrm D}(w),\qquad
 a_D(w)=3^{-|w|}(3/5)^{N_{\mathrm D}(w)}.
\]
This identity places the entire revisit factor in the coefficient.
Partition-function ratios telescope to
\(R_D^\bullet(w)=\prod_{j=1}^{|w|}\rho_j^\bullet\), where
\(\rho_j^\bullet\) is the ratio for the \(j\)-th step.

\begin{lemma}\label{l:stepid}
For vectors \(a,u_1,\ldots,u_k\in\mathbb R^3\),
\[
 \mathcal M\!\left((a\cdot n)\prod_{i=1}^k(1+n\cdot u_i)\right)
 =\frac13\sum_{i=1}^k(a\cdot u_i)
       \lambda_{\mathrm D}\!\left(\prod_{i'\ne i}(1+n\cdot u_{i'})\right)
       \quad(0\le k\le3),
\]
and
\[
 \lambda_{\mathrm D}\!\left((a\cdot n)\prod_{i=1}^k(1+n\cdot u_i)\right)
       =\frac15\sum_{i=1}^k a\cdot u_i\quad(0\le k\le1).
\]
An empty sum is zero. The second identity removes the departure site
with no residual scalar mass.
\end{lemma}

\begin{proof}
Expand the product into subsets of \(\{1,\ldots,k\}\).
In the first identity, the inserted linear factor permits only
subsets of sizes one and three to contribute, by
Lemma~\ref{l:moments}. The value is
\[
 \frac13\sum_i a\cdot u_i+
 \frac1{15}\sum_{i<j<l}
 \bigl((a\cdot u_i)(u_j\cdot u_l)
       +(a\cdot u_j)(u_i\cdot u_l)
       +(a\cdot u_l)(u_i\cdot u_j)\bigr).
\]
For the residual product, which has at most two factors,
\[
 \lambda_{\mathrm D}\left(\prod_{i'\ne i}(1+n\cdot u_{i'})\right)
 =1+\frac15\sum_{\substack{j<l\\j,l\ne i}}u_j\cdot u_l.
\]
Multiplying by \((a\cdot u_i)/3\) and summing over \(i\)
gives exactly the preceding expression, including the empty sums
when \(k=0,1,2\).

For the second identity, \(k=0\) gives zero by the first moment.
For \(k=1\), the constant part again vanishes and the second
moment is
\((3/5)(a\cdot u_1)/3=(a\cdot u_1)/5\).
There is no residual factor depending on \(n\).
Thus the site can be deleted with the displayed coefficient and
with no change of the normalization of the remaining site laws.
All polynomials evaluated here have degree at most four.
\end{proof}

\begin{proposition}\label{p:pairwalk}
Let \(D\) be a domain or auxiliary network, and let \(u\) be an
admissible initial start. Let \(v\ne u\) be nonremoved; if it is
\(\mathrm U\), assume at most three incident factors, counting
original legs; if it is \(\mathrm D\), assume at most one present bond.
Put \(T=S\cup\{v\}\). For \(\bullet=\tau,0\),
\[
 \mathbb E_D^\bullet(n_u\cdot n_v)=
 \sum_{w\in\mathcal T_D(u,T)}a_D(w)R_D^\bullet(w)
           \mathbb E_{D[w]}^\bullet(n_{v_{|w|}}\cdot n_v).
\]
If the endpoint is \(v\), its last factor is independent of the
choice of \(\bullet\): it is \(1\) for status \(\mathrm U\),
and \(3/5\) for status \(\mathrm D\).
Every partition product with an inserted observable used in this
identity has separate degree at most four.
\end{proposition}

\begin{proof}
We prove an identity of unnormalized numerators. Write
\(n_u\cdot n_v=\sum_{\beta=1}^3n_{u,\beta}n_{v,\beta}\), and fix
\(\beta\). At any nonterminal prefix ending at \(y\), the marked
factor is \(n_{y,\beta}n_{v,\beta}\), with \(y\ne v\) and
\(y\notin S\). Sum over \(n_y\) first.
If its status is \(\mathrm U\), the first identity of
Lemma~\ref{l:stepid}, with \(a=e_\beta\) and the present
neighbouring spins as the \(u_i\), gives
\[
 Z^\bullet_E[n_{y,\beta}n_{v,\beta}]
   =\frac13\sum_{o:\,\{y,o\}\text{ present in }E}
        Z^\bullet_{E[(y,o)]}[n_{o,\beta}n_{v,\beta}].
\]
Here \(E\) denotes the current network; \(E[(y,o)]\) is its
one-step trace, so its law at \(y\) is
\(\lambda_{\mathrm D}\) and the selected bond has been dropped.
If \(y\) is \(\mathrm D\), the second identity gives instead
\[
 Z^\bullet_E[n_{y,\beta}n_{v,\beta}]
   =\frac15\sum_{o:\,\{y,o\}\text{ present in }E}
        Z^\bullet_{E[(y,o)]}[n_{o,\beta}n_{v,\beta}],
\]
where \(y\) is removed in every term.
The departure bounds needed for these identities were verified
after Definition~\ref{d:trace}. There are no leg factors at \(y\).
If there is no present bond, either displayed sum is zero.

Iterate these equalities until the marked site reaches \(T\).
Every iteration consumes a previously unused present unoriented
bond. There are only finitely many bonds initially, so the
iteration terminates after at most that many steps, even for
branches which end with an empty sum. No unused present bond is
destroyed by the trace, as proved after
Definition~\ref{d:trace}. Hence the branches that reach \(T\) are
in bijection with the trails in \(\mathcal T_D(u,T)\): a sequence
of choices determines one such trail and conversely. The product of its
coefficients is \(a_D(w)\). We have proved
\[
 Z^\bullet_D[n_{u,\beta}n_{v,\beta}]
 =\sum_{w\in\mathcal T_D(u,T)}
       a_D(w)Z^\bullet_{D[w]}
                   [n_{v_{|w|},\beta}n_{v,\beta}].
\]
Sum over \(\beta\) and divide once by the positive \(Z^\bullet_D\).
This proves the expansion.

If the endpoint is \(v\), this is its first hit, so the status of
\(v\) has not changed. For status \(\mathrm U\), its spin has
squared norm one pointwise. For status \(\mathrm D\), its sole
possible present bond in \(D\) must have been the arrival bond;
there are no original legs at a \(\mathrm D\) site. After arrival
it is isolated and has its normalized law
\(\lambda_{\mathrm D}\), independently of all other spins.
Thus its expected squared norm is \(3/5\), for either pin choice.

For completeness, before a terminal hit the fixed target \(v\)
has at most three factors if \(\mathrm U\), or at most one
present bond if \(\mathrm D\). Inserting its one coordinate
therefore gives degree at most four or two, respectively.
The moving marked site has at most three bonds if \(\mathrm U\),
or one if \(\mathrm D\), and no legs; its inserted coordinate
has the same respective bounds. At a terminal site in
\(S\setminus\{v\}\), the lost arrival bond leaves at most three
factors, counting all unchanged original legs, and the inserted
coordinate gives degree at most four. At a terminal \(v\) of
status \(\mathrm U\), at most two factors remain after arrival;
inserting \(|n_v|^2\) again gives degree at most four.
A terminal \(v\) of status \(\mathrm D\) is isolated and has
only the degree-two insertion. All unmarked sites have at most
four factors, and changing a law or deleting a bond raises no
degree. These cases cover every numerator in the iteration.
\end{proof}

\begin{definition}\label{d:ratios}
At a step \(j\), put
\(\rho_j^\bullet=Z^\bullet_{D[w_j]}/Z^\bullet_{D[w_{j-1}]}\).
At a \(\mathrm D\) departure \(\rho_j^\tau=\rho_j^0=1\).
At a \(\mathrm U\) departure from \(y\) to \(o\), let \(E'\)
be the pre-step network with \(y\) removed.
If the other present neighbours are \(f,g\), set
\[
 I=1+\tfrac13(P_{fg}+P_{fo}+P_{go}),\quad
 J=1+\tfrac15P_{fg},\quad P_{ab}=n_a\cdot n_b,\qquad
 (\rho_j^\bullet)^{-1}=\mathbb E_{E'}^\bullet I/
                                      \mathbb E_{E'}^\bullet J.
\]
With only one other neighbour \(f\), use \(I=1+P_{fo}/3,\ J=1\);
with no other neighbour, \(\rho_j^\bullet=1\).
Straight/turning refers to the incoming direction, with a designated
absent incoming bond at the first step.
\end{definition}

These formulas follow by summing over the departure spin before
performing the remaining product sum. With three neighbours, the
pre-step \(\lambda_{\mathrm U}\) sum is \(I\), by its second
moment and vanishing odd moments. The post-step
\(\lambda_{\mathrm D}\) sum over its two retained bonds is \(J\),
by its second moment \(1/5\). Thus
\[
 Z^\bullet_{D[w_{j-1}]}=Z^\bullet_{E'}\mathbb E^\bullet_{E'}I,
 \qquad
 Z^\bullet_{D[w_j]}=Z^\bullet_{E'}\mathbb E^\bullet_{E'}J.
\]
With two neighbours these same sums are \(1+P_{fo}/3\) and \(1\);
with one they are both \(1\). At a \(\mathrm D\) departure
there is at most the selected bond, and
\(\lambda_{\mathrm D}(1+n\cdot n_o)=1\), so deletion leaves
the partition function unchanged. All divisions are justified by
Lemma~\ref{l:finitefacts}. These computations also show directly
that the revisit has ratio one for both pin choices.

Incoming directions are lattice unit vectors, namely
\((\pm1,0)\) or \((0,\pm1)\).
At an admissible start, all four neighbouring sites lie in \(G\)
because the start is outside \(S\); there are at most three
present bonds. Hence an absent bond can always be designated as
the initial incoming bond. At every later step the actual incoming
bond has already been deleted at both ends.

\begin{lemma}\label{l:response}
At a \(\mathrm U\) departure \(y\notin S\), write
\(r_{ab}=\mathbb E^\tau_{E'}P_{ab}-\mathbb E^0_{E'}P_{ab}\).
With three present neighbours,
\[
 \log\rho_j^0-\log\rho_j^\tau
       =\alpha_jr_{fg}+\beta_j(r_{fo}+r_{go}),\qquad
 0<\alpha_j\le\bar\alpha=\frac{18094}{100000},\quad
 0<\beta_j\le\bar\beta=\frac{36751}{100000}.
\]
With two present neighbours the response is \(\beta_j'r_{fo}\),
where \(0<\beta_j'\le35510/100000\le\bar\beta\);
with one it is zero. Also \(\rho_j^0\le1\) and
\(\rho_j^\tau\le12/5\).
A \(\mathrm D\) departure has zero logarithmic response.
\end{lemma}

\begin{proof}
For positive \(x,y\), define
\[
 L(x,y)=\int_0^1x^t y^{1-t}\,\mathrm dt
 =\begin{cases}
 (x-y)/(\log x-\log y),&x\ne y,\\
 x,&x=y.
 \end{cases}
\]
The equality follows by integrating the exponential when \(x\ne y\).
The integral shows that \(L\) is positive, increasing in each argument,
and between the minimum and maximum of its arguments. In both cases
\(\log x-\log y=(x-y)/L(x,y)\).

First suppose that the present neighbours are \(f,g,o\).
Write \(I^\bullet=\mathbb E^\bullet_{E'}I\),
\(J^\bullet=\mathbb E^\bullet_{E'}J\), and
\(L_I=L(I^\tau,I^0)\), \(L_J=L(J^\tau,J^0)\).
Using the ratio formula in Definition~\ref{d:ratios} gives
\[
 \begin{split}
 \log\rho_j^0-\log\rho_j^\tau
 &=\frac{I^\tau-I^0}{L_I}
       -\frac{J^\tau-J^0}{L_J}\\
 &=\left(\frac1{3L_I}-\frac1{5L_J}\right)r_{fg}
       +\frac1{3L_I}(r_{fo}+r_{go}).
 \end{split}
\]
We use these displayed expressions to define \(\alpha_j,\beta_j\),
also when a response or both arguments of a logarithmic mean coincide.

Since \(y\notin S(G)\), it has distance at least three from
\(\mathbb Z^2\setminus G\). Each of its neighbours has distance
at least two from that complement and carries no original leg.
After \(y\) is removed, each has at most three present bonds.
The neighbours are pairwise nonadjacent. On every positive fibre
of the other spins in \(E'\), Lemmas~\ref{l:finitefacts}
and~\ref{l:35} therefore give independent conditional spins with
means \(m_f,m_g,m_o\) of norms at most \(3/5\), under either
pin choice. Put \(Q=\sum_{\{a,b\}\subset\{f,g,o\}}m_a\cdot m_b\).
Then
\[
 -\frac{27}{50}\le
 \frac{|\sum_a m_a|^2-\sum_a|m_a|^2}{2}=Q
 \le\sum_{\{a,b\}}|m_a||m_b|\le\frac{27}{25}.
\]
Also \(|m_f\cdot m_g|\le9/25\).
Averaging these conditional inequalities proves
\[
 \frac{41}{50}\le I^\bullet\le\frac{34}{25},\qquad
 \frac{116}{125}\le J^\bullet\le\frac{134}{125}.
\]
Every pair insertion used here is admissible for free positivity:
removal of \(y\) leaves at most three factors at each inserted
endpoint, and the pair adds one degree there.
Thus Lemma~\ref{l:finitefacts} gives \(I^0\ge1\), \(J^0\ge1\).
The monotonicity of \(L\) now gives
\[
 L_I\ge L(41/50,1),\quad L_I\le34/25,\qquad
 116/125\le L_J\le134/125.
\]
In particular
\[
 \alpha_j\ge\frac{25}{102}-\frac{25}{116}
       =\frac{175}{5916}>0,\qquad \beta_j>0.
\]

Here are exact rational upper bounds for the logarithms needed in
the upper estimates. For \(0<s<1\), set \(t=(1-s)/(1+s)\).
The substitution giving
\(-\log s=2\int_0^t(1-u^2)^{-1}\,\mathrm du\), together with
\[
 \frac1{1-u^2}=1+u^2+\frac{u^4}{1-u^2},\qquad
 \frac1{1-u^2}\le\frac1{1-t^2}\quad(0\le u\le t),
\]
proves
\[
 \frac{-\log s}{1-s}\le
 B(s):=\frac2{1-s}
       \left(t+\frac{t^3}{3}+\frac{t^5}{5(1-t^2)}\right).
\]
Substitution and rational simplification give
\[
 \begin{array}{c|c|c}
 s&B(s)/3&\text{rational upper bound}\\ \hline
 41/50&340634561/926892330&367502/10^6\\
 22/25&24332081/68523180&35510/100000
 \end{array}
\]
The middle entries are strictly smaller than the right entries;
all denominators are positive. For example their cross-product
differences are respectively \(224059660\) and \(50021800\).
Consequently
\[
 \beta_j\le\frac1{3L(41/50,1)}
       \le\frac{367502}{10^6}<\bar\beta,\qquad
 \alpha_j\le\frac{367502}{10^6}-\frac{25}{134}
       =\frac{6061317}{33500000}<\bar\alpha.
\]
The last strict inequality has difference
\(173/33500000>0\).

With just the neighbours \(f,o\), \(J=1\) and
\[
 I^\bullet\in[22/25,28/25],\qquad I^0\ge1.
\]
The same logarithmic identity gives
\(\beta_j'=1/(3L(I^\tau,I^0))>0\) and response
\(\beta_j'r_{fo}\). The second row of the table proves
\(\beta_j'\le35510/100000\le\bar\beta\).
With only \(o\), both partition ratios are one and the response
is zero.

For the free ratio with three neighbours, free positivity also gives
\[
 I^0-J^0
 =\frac2{15}\mathbb E^0_{E'}P_{fg}
      +\frac13\mathbb E^0_{E'}(P_{fo}+P_{go})\ge0.
\]
For two neighbours use \(I^0-1=\mathbb E^0_{E'}P_{fo}/3\ge0\);
for one neighbour the ratio is one. Hence \(\rho_j^0\le1\)
in all cases. Pointwise, for three neighbours of norm at most one,
\[
 \sum_{\{a,b\}}n_a\cdot n_b
   =\frac{|\sum_a n_a|^2-\sum_a|n_a|^2}{2}\ge-\frac32,
 \qquad I\ge\frac12,\qquad J\le\frac65.
\]
It follows that \(\rho_j^\tau=J^\tau/I^\tau\le12/5\).
For two neighbours, \(I\ge2/3\), \(J=1\), so the bound is
\(3/2\le12/5\); for one it is one.
Finally a \(\mathrm D\) departure has both ratios equal to one
by Definition~\ref{d:ratios}, proving its zero logarithmic response.
\end{proof}

\begin{lemma}\label{l:exposure}
For an admissibly started stopped trail, let \(\Phi\) be its
\(\mathrm U\) departures, and suppose \(\rho_j^\tau\le x_j\),
\(x_j\ge1\), for \(j\in\Phi\).
For every ordering \(\sigma\) of \(\Phi\),
\[
 |R_D^\tau(w)-R_D^0(w)|
 \le\sum_{k=1}^{|\Phi|}
       |\log\rho_{\sigma(k)}^0-\log\rho_{\sigma(k)}^\tau|
                         \prod_{l\le k}x_{\sigma(l)}.
\]
Multiplication by \(a_D(w)\) gives the corresponding bound on the
coefficient in the pair expansion.
\end{lemma}

\begin{proof}
All step ratios are positive. The \(\mathrm D\) ratios are one
for both pin choices, so
\[
 R_D^\tau(w)=\prod_{j\in\Phi}b_j,\qquad
 R_D^0(w)=\prod_{j\in\Phi}c_j,\qquad
 b_j=\rho_j^\tau,\quad c_j=\rho_j^0.
\]
For \(m=|\Phi|\), telescoping in the given order gives the exact
identity
\[
 \prod_{j\in\Phi}b_j-\prod_{j\in\Phi}c_j
 =\sum_{k=1}^m
   \left(\prod_{l<k}b_{\sigma(l)}\right)
   (b_{\sigma(k)}-c_{\sigma(k)})
   \left(\prod_{l>k}c_{\sigma(l)}\right).
\]
For positive \(b,c\), integration of the derivative of
\(\exp t\) between \(\log b\) and \(\log c\) yields
\[
 |b-c|\le\max\{b,c\}\,|\log b-\log c|.
\]
Lemma~\ref{l:response} gives \(c_j\le1\), and the hypothesis
gives \(b_j\le x_j\), with \(x_j\ge1\).
Thus \(\max\{b_j,c_j\}\le x_j\), and the absolute value of
the \(k\)-th telescoping term is at most
\[
 |\log\rho_{\sigma(k)}^0-\log\rho_{\sigma(k)}^\tau|
          \prod_{l\le k}x_{\sigma(l)}.
\]
Sum these inequalities. If \(m=0\), both products are one and
the sum is empty, so the same assertion holds.
Finally \(a_D(w)>0\), and multiplying the inequality by it
gives the stated coefficient bound.
\end{proof}

\begin{definition}\label{d:pattern}
Let \(\mathsf T=\{z:d(z,0)\le2\}\), with passage site \(0\),
arrival site \(\mathsf s=(0,-1)\), other neighbours
\(\mathsf e=(1,0),\mathsf N=(0,1),\mathsf W=(-1,0)\),
four corners \((\pm1,\pm1)\), and tips \((\pm2,0),(0,\pm2)\).
The external slots are the two external bonds at each corner and
the three at each tip, twenty in total.
A pattern chooses an out-neighbour \(o\in\{\mathsf e,\mathsf N,\mathsf W\}\)
and statuses/kept sets with the following rules.
The passage site is \(\mathrm U\) and keeps precisely its three
non-arrival bonds. The arrival site is removed or is \(\mathrm D\)
keeping two of its three non-passage bonds.
Each other neighbour is \(\mathrm U,\mathrm X\), or \(\mathrm D\)
keeping its passage bond and one other bond; \(o\) is not removed.
Each corner is \(\mathrm U,\mathrm X\), or \(\mathrm D\) keeping
two of its four bonds. Every \(\mathrm U\) site other than the
passage site keeps all four bonds.
Between nonremoved non-tip sites, a bond is kept by both ends or
by neither. Tips are covering \(\mathrm U\) sites; their bond to
a neighbour \(z\) is present exactly when \(z\) is nonremoved and
keeps that bond. No consistency condition with a global tip kept
set is imposed.
For \(m\in(\mathbb S^2)^{20}\), include every external slot factor
at every nonremoved corner and tip, including both corner slots
when the corresponding actual external bonds are missing.
Let \(Z_{\rm cur}(\pi;m)\) be the finite product functional of
all present template bonds and these slots, using the status laws.
Let \(Z_{\rm next}(\pi;m)\) change the passage law to
\(\lambda_{\mathrm D}\) and drop its out-bond.
The resulting class is \(\mathcal P_{\rm ext}\); it is straight
when \(o=\mathsf N\), and turning otherwise.
\end{definition}

Both local sums are strictly positive for every unit slot choice.
Indeed a corner has at most two internal factors and two slots,
a tip has at most one internal factor and three slots, a neighbour
has at most its four internal factors, and the passage site has
at most three before the step and two after it. Removed sites
contribute no variable or slot. Lemma~\ref{l:finitefacts}
therefore applies to each sum. Every slot vector occurs once,
affinely; this property also applies when a slot at a removed site
has no effect.

\begin{definition}[(H1)]\label{h:H1}
For \(c_{\mathrm s},c_{\mathrm t}\in(0,1]\), (H1) means that
for every \(\pi\in\mathcal P_{\rm ext}\) and every
\(m\in(\mathbb S^2)^{20}\),
\[
 Z_{\rm cur}(\pi;m)\ge c_{\rm type(\pi)}Z_{\rm next}(\pi;m).
\]
In the paper the fixed thresholds are
\(c_{\mathrm s}=961/1000\), \(c_{\mathrm t}=187/200\).
No minimum, minimizing pins, or optimal inversion is asserted.
\end{definition}

We use (H1) through its pointwise inequality on the full covered
pattern class of Definition~\ref{d:pattern}.

\begin{lemma}\label{l:local}
At any \(\mathrm U\) departure outside \(S\) whose normalized
non-tip template data belong to \(\mathcal P_{\rm ext}\), (H1)
implies \(\rho_j^\tau\le1/c_{\rm type}\).
The conclusion is uniform in all original unit pins, all outside
spins of norm at most one, all missing external factors, and all
actual tip statuses and kept sets allowed by the ordinary-domain rules.
\end{lemma}

\begin{proof}
Let \(E\) and \(E^+\) be the pre-step and post-step networks,
and let \(\mathsf T(y)=y+\mathsf T\) be the template.
Because \(y\notin S(G)\), this template lies in \(G\).
Apply a lattice symmetry to put the incoming direction and the
out-neighbour in the positions of Definition~\ref{d:pattern}.
This just relabels the sites and their spins.

Fix the outside spins at \(\omega\). Split off all site weights
and factors involving only sites outside the template; their
product, denoted \(H(\omega)\), is nonnegative and identical
for \(E\) and \(E^+\). Let \(C(\omega)\), \(N(\omega)\) be the
remaining template sums for \(E\), \(E^+\), respectively.
There is a factor \(1+n_z\cdot b\) for each actual external
bond or original leg of a nonremoved template site \(z\), with
\(|b|\le1\). Original pins have \(|b|=1\).
Each external direction carries at most one such factor.
No external factor is incident to the passage site or its
neighbours, since these sites are at distances at least three
and two, respectively, from the original complement.

We first replace the actual tips. A tip \(t\) has precisely one
possible neighbour \(z\) inside the template; its other three
directions are external slots. If \(t\) is nonremoved, write
\[
 W_t(n)=\prod_{\text{actual external factors at }t}(1+n\cdot b),
 \quad A_t=\lambda_{\mathrm{st}(t)}(W_t)>0,\quad
 b_t=\frac{\lambda_{\mathrm{st}(t)}(nW_t)}{A_t}.
\]
There are at most three external factors, so \(A_t>0\) and
\(|b_t|\le3/5\), by the proof of Lemma~\ref{l:35}.
If its internal bond is present, summing over this tip gives
\(A_t(1+b_t\cdot n_z)\); if the bond is absent it gives \(A_t\).
A removed tip contributes the constant one.
These sums and constants are identical in the current and next
template, since the step changes only the passage site and its
out-bond.

In the covering pattern the bond from this tip to \(z\) is
required exactly when \(z\) is nonremoved and keeps it.
If that bond is required but the actual tip has no present bond,
use the desired effective mean zero and the actual constant
as \(A_t\), taking \(A_t=1\) for a removed tip.
If it is required and actually present, use the effective mean
\(b_t\) just computed. In either case Lemma~\ref{l:35}
supplies three unit external fields for a covering
\(\mathrm U\) tip with that exact mean. Let \(A'_t>0\) be
its external three-factor normalizer. Its integrated factor
is the desired factor multiplied by \(A'_t/A_t\).
If the covering bond is not required, the actual bond is
absent as well; choose any three unit fields for the covering
tip and use the same positive normalizer adjustment.
Thus replacing all four tips gives a covered-tip template whose
current and next sums are multiplied by the same strictly
positive scalar \(\prod_t A'_t/A_t\).
The zero mean and the full ball of means used here are attained
by the explicit moment formula in Lemma~\ref{l:35}.

We next fill and resolve the corner slots. For each external
direction at a nonremoved corner, use its actual field
\(b\) if present, and use \(b=0\) if absent.
For any \(|b|\le1\), choose a unit \(e\) parallel to \(b\)
when \(b\ne0\), and any unit \(e\) when \(b=0\). Then
\[
 1+n_z\cdot b
 =\frac{1+|b|}{2}(1+n_z\cdot e)
       +\frac{1-|b|}{2}(1-n_z\cdot e).
\]
This is also the filling identity for an absent factor.
Expanding independently over the at most eight corner slots
expresses the current and next sums as convex combinations
with identical coefficients of sums with unit vectors in every
corner slot. Both slots are filled at every nonremoved corner,
regardless of its actual external kept bonds.
At removed corners the slot values can be assigned arbitrarily.

By the premise, the non-tip statuses and internal bonds are
exactly those of a pattern \(\pi\in\mathcal P_{\mathrm{ext}}\).
After the two replacements, every term is therefore
\(Z_{\rm cur}(\pi;m)\) or \(Z_{\rm next}(\pi;m)\) for the same
unit slot assignment \(m\); its type is the type of the given
step. Apply (H1) term by term, average with the nonnegative
corner coefficients, and undo the common positive tip scalar.
This proves
\(C(\omega)\ge c_{\rm type}N(\omega)\)
for every outside configuration \(\omega\).
Multiply by \(H(\omega)\) and sum over the outside spins.
The result is
\[
 Z_E^\tau\ge c_{\rm type}Z_{E^+}^\tau.
\]
Both partition functions are positive by
Lemma~\ref{l:finitefacts}, so division proves
\(\rho_j^\tau=Z_{E^+}^\tau/Z_E^\tau\le1/c_{\rm type}\).
The construction allowed each fixed outside field to range
over the whole closed unit ball, and every missing slot and
allowed tip status was treated explicitly. This proves all
the asserted uniformity.
\end{proof}

\begin{definition}\label{d:weights}
Put \(x_{\mathrm s}=1/c_{\mathrm s}\), \(x_{\mathrm t}=1/c_{\mathrm t}\).
A lattice trail \(w=(v_0,\ldots,v_n)\) has distinct unoriented
nearest-neighbour bonds. Give it a designated unit incoming direction
\(a\), and forbid \(v_1-v_0=-a\) when \(n\ge1\); the virtual
incoming bond is not itself a used bond of the trail.
For departure \(i\), \(0\le i<n\), give weight \(1/5\) if \(v_i\)
has been departed from earlier, and otherwise weight \(x_{\mathrm s}/3\)
for continuing in the incoming direction and \(x_{\mathrm t}/3\)
for a turn. Let \(\mathsf W(w)\) be the product, with
\(\mathsf W(\varnothing)=1\).
\end{definition}

At a later departure the incoming direction is
\(v_i-v_{i-1}\); at \(i=0\) it is \(a\).
Distinctness of the unoriented bonds forbids immediate reversal
at every later step. The first-step rule is an additional
restriction specified by the incoming direction.
In a trace, a site which has been departed from earlier cannot
again have status \(\mathrm U\). At any subsequent departure
from that site it is \(\mathrm D\), with coefficient \(1/5\)
and ratio one. A site initially of status \(\mathrm D\) may
instead have its first lattice departure weighted as a first
visit; since \(c_{\rm type}\le1\),
\[
 \frac15\le\frac13\le\frac{x_{\rm type}}3.
\]
This is the comparison for initially \(\mathrm D\) sites.

\begin{definition}[(H2)]\label{h:H2}
For \(C_V\ge1\) and \(0<q<1\), (H2) means that for every
integer \(n\ge0\), every start and every designated incoming direction,
\[
 \sum_{\substack{|w|=n\\\text{initial reversal forbidden}}}
                      \mathsf W(w)\le C_Vq^n .
\]
The sum is over plane lattice trails with those specified initial data.
\end{definition}


\begin{lemma}\label{l:onepoint}
Assume (H1) and (H2). Start a trace at an admissible \(x\notin S\)
with incoming direction \(a\) such that \(\{x-a,x\}\) is absent.
Suppose each stopped trail from \(x\) to \(S\) has at most twelve
\(\mathrm U\) departures whose templates are not covered patterns.
Then
\[
 \sum_{w\in\mathcal T_D(x,S)}a_D(w)
                  \max\{R_D^\tau(w),R_D^0(w)\}
 \le (12/5)^{12}\frac{C_V}{1-q}\,q^{d(x,S)}.
\]
Pointwise the summand is at most \((12/5)^{12}\mathsf W(w)\).
\end{lemma}

\begin{proof}
The absent bond \(\{x-a,x\}\) ensures that every trail of the
network from \(x\) has the forbidden initial reversal required
in Definition~\ref{d:weights}. Subsequent steps have distinct
unoriented bonds, so every stopped trail in the left-hand side
is among the plane trails in (H2) with these same initial data.

Fix such a stopped trail \(w\), and let \(\omega_j\) be its
\(j\)-th lattice weight in \(\mathsf W(w)\).
We bound its coefficient for both choices of \(\bullet\)
by the same product of step majorants.
At a \(\mathrm D\) departure the product of the departure
coefficient and \(\rho_j^\bullet\) is exactly \(1/5\).
If this site was departed from earlier in \(w\), this is
exactly \(\omega_j\). If it was not, then
\(\omega_j=x_{\rm type}/3\ge1/3\), so it still bounds \(1/5\).
This includes an initially \(\mathrm D\) start, whose
admissibility ensures at most one present bond.

A \(\mathrm U\) departure is necessarily the first departure
from its site in \(w\). If its pattern is covered,
Lemma~\ref{l:local} gives \(\rho_j^\tau\le x_{\rm type}\)
and Lemma~\ref{l:response} gives
\(\rho_j^0\le1\le x_{\rm type}\). Thus, for both choices,
\[
 \frac13\rho_j^\bullet\le\frac{x_{\rm type}}3=\omega_j.
\]
At an uncovered \(\mathrm U\) departure, the same uniform
response lemma gives \(\rho_j^\tau\le12/5\),
\(\rho_j^0\le1\le12/5\). Therefore
\[
 \frac13\rho_j^\bullet\le\frac45
          \le(12/5)\,\omega_j,
\]
because \(\omega_j\ge1/3\).
There are at most twelve such departures, wherever they occur
along \(w\). Multiplying the step estimates and telescoping
the ratios proves, for both pin choices with the same bound,
\[
 a_D(w)R_D^\bullet(w)
   \le(12/5)^{12}\prod_j\omega_j
   =(12/5)^{12}\mathsf W(w).
\]
Taking the maximum preserves this pointwise inequality.

The length of a nearest-neighbour trail from \(x\) to \(S\)
is at least \(d_0=d(x,S)\). If there are no stopped trails,
the asserted sum is zero. Otherwise, group the finite domain
sum by length and enlarge each group to all plane trails
of that length with the fixed start and incoming direction.
Using (H2), nonnegativity and \(0<q<1\), we obtain
\[
 \begin{split}
 \sum_{w\in\mathcal T_D(x,S)}
       a_D(w)\max\{R_D^\tau(w),R_D^0(w)\}
 &\le(12/5)^{12}\sum_{n\ge d_0}
              \sum_{\substack{|w|=n\\
                    \text{initial reversal forbidden}}}\mathsf W(w)\\
 &\le(12/5)^{12}C_V\sum_{n\ge d_0}q^n\\
 &=(12/5)^{12}\frac{C_V}{1-q}\,q^{d_0}.
 \end{split}
\]
Here \(d_0\) is a finite integer: the existence of \(x\) makes
\(G\) nonempty, and every nonempty finite \(G\subset\mathbb Z^2\)
has a site adjacent to its complement, hence \(S\ne\varnothing\).
This completes the proof.
\end{proof}

